% Part: normal-modal-logic
% Chapter: completeness
% Section: complete-consistent-sets

\documentclass[../../../include/open-logic-section]{subfiles}

\begin{document}

\olfileid{nml}{com}{ccs}

\olsection{Himpunan $\Sigma$-Konsisten Jangkep}

Upamane $\Sigma$ himpunan !!{formula}s modal---anggepen minangka
aksioma utawa prinsip kang nemtokake logika modal normal. Himpunan $\Gamma$
iku $\Sigma$-konsisten yen lan mung yen $\Gamma \Proves/[\Sigma] \lfalse$,
yaiku yen ora ana !!{derivation} saka~$!A_1 \lif (!A_2 \lif \cdots (!A_n \lif
\lfalse)\dots)$ saka $\Sigma$, kanthi saben $!A_i \in \Gamma$. Kita bakal
mbangun model ``kanonik'' kang saben jagade dianggep minangka himpunan
$\Sigma$-konsisten kanthi sipat mligi: ora mung $\Sigma$-konsisten,
nanging uga maksimal ing teges bisa nemtokake nilai bebener saben
!!{formula} modal: kanggo saben $!A$, $!A
\in \Gamma$ utawa $\lnot !A \in \Gamma$:

\begin{defn}
  Himpunan $\Gamma$ diarani \emph{$\Sigma$-konsisten jangkep} yen lan
  mung yen himpunan iku $\Sigma$-konsisten lan kanggo saben~$!A$,
  $!A \in \Gamma$ utawa $\lnot !A \in \Gamma$.
\end{defn}

Himpunan $\Sigma$-konsisten jangkep $\Gamma$ nduweni sawetara
sipat kang migunani. Kapisan, himpunan iku katutup sacara deduktif,
yaiku yen $\Gamma \Proves[\Sigma] !A$, mula $!A \in \Gamma$.
Iki mligine ateges saben instansi saka !!a{formula}~$!A \in \Sigma$
uga $\in \Gamma$. Kajaba iku, keanggotaan ing $\Gamma$ nggambarake
sarat bebener konektif proposisional. Iki bakal wigati nalika kita
netepake ``model kanonik.''

\begin{prop}\ollabel{prop:ccs-properties}
  Upamane $\Gamma$ $\Sigma$-konsisten jangkep. Banjur:
  \begin{enumerate}
  \item \ollabel{prop:ccs-closed}%
    $\Gamma$ katutup sacara deduktif ing~$\Sigma$.
  \item \ollabel{prop:ccs-sigma}%
    $\Sigma \subseteq \Gamma$.
  \tagitem{prvFalse}{\ollabel{prop:ccs-lfalse}%
    $\lfalse \notin \Gamma$}{}
  \tagitem{prvTrue}{\ollabel{prop:ccs-ltrue}%
    $\ltrue \in \Gamma$}{}
  \item \ollabel{prop:ccs-lnot}%
    $\lnot!A \in \Gamma$ yen lan mung yen $!A \notin
    \Gamma$.
  \tagitem{prvAnd}{\ollabel{prop:ccs-land}%
    $!A \land !B \in \Gamma$ yen lan mung yen $!A \in \Gamma$ lan $!B \in \Gamma$}{}
  \tagitem{prvOr}{\ollabel{prop:ccs-lor}%
    $!A \lor !B \in \Gamma$ yen lan mung yen $!A \in \Gamma$ utawa $!B \in \Gamma$}{}
  \tagitem{prvIf}{\ollabel{prop:ccs-lif}%
    $!A \lif !B \in \Gamma$ yen lan mung yen $!A \notin \Gamma$ utawa $!B \in \Gamma$}{}
  \tagitem{prvIff}{\ollabel{prop:ccs-liff}%
    $!A \liff !B \in \Gamma$ yen lan mung yen $!A \in \Gamma$ lan $!B \in
    \Gamma$, utawa $!A \notin \Gamma$ lan $!B \notin \Gamma$}{}
  \end{enumerate}
\end{prop}

\begin{proof}
  \begin{enumerate}
  \item Upamane $\Gamma \Proves[\Sigma] !A$ nanging $!A \notin
    \Gamma$. Amarga $\Gamma$ $\Sigma$-konsisten jangkep,
    $\lnot!A \in \Gamma$. Iki bakal ndadekake $\Gamma$ ora konsisten,
    amarga $!A, \lnot !A \Proves[\Sigma] \lfalse$.

\item Yen $!A \in \Sigma$, mula $\Gamma \Proves[\Sigma] !A$, lan $!A
    \in \Gamma$ amarga katutup sacara deduktif, yaiku
    kasus~\olref{prop:ccs-closed}.

\tagitem{prvFalse}{Yen $\lfalse \in \Gamma$, mula $\Gamma
    \Proves[\Sigma] \lfalse$, saengga $\Gamma$ bakal ora
    $\Sigma$-konsisten.}{}

\tagitem{prvTrue}{$\Gamma \Proves[\Sigma] \ltrue$, mula $\ltrue \in
    \Gamma$ amarga katutup sacara deduktif, yaiku
    kasus~\olref{prop:ccs-closed}.}{}

\item Yen $\lnot!A \in \Gamma$, mula amarga konsistensi $!A \notin
    \Gamma$; lan yen $!A \notin \Gamma$, mula $\lnot !A \in \Gamma$
    amarga $\Gamma$ $\Sigma$-konsisten jangkep.

\tagitem{prvAnd}{\iftag{probAnd}{Latihan.}{Upamane $!A \land !B \in
      \Gamma$. Amarga $(!A \land !B) \lif !A$ iku instansi tautologis,
      $!A \in \Gamma$ amarga katutup sacara deduktif, yaiku
      kasus~\olref{prop:ccs-closed}. Kanthi cara padha, $!B \in \Gamma$.
      Kosok baline, upamane $!A \in \Gamma$ lan $!B \in
      \Gamma$. Banjur katutup sacara deduktif njalari $(!A \land !B) \in
      \Gamma$, amarga $!A \lif (!B \lif (!A \land !B))$ iku
      instansi tautologis.}}{}

\tagitem{prvOr}{\iftag{probOr}{Latihan.}{Upamane $!A \lor !B \in
      \Gamma$, nanging $!A \notin \Gamma$ lan $!B \notin \Gamma$.
      Amarga $\Gamma$ $\Sigma$-konsisten jangkep, $\lnot!A \in \Gamma$
      lan $\lnot !B \in \Gamma$. Banjur $\lnot(!A \lor !B) \in \Gamma$
      amarga $\lnot !A \lif (\lnot !B \lif \lnot (!A \lor !B))$ iku
      instansi tautologis. Iki ndadekake $\Gamma$ ora
      $\Sigma$-konsisten, sawijining kontradiksi. Kosok baline, yen
      $!A \in \Gamma$ utawa $!B \in \Gamma$, mula $!A \lor !B \in \Gamma$
      amarga katutup sacara deduktif, awit $!A \lif (!A \lor !B)$ lan
      $!B \lif (!A \lor !B)$ iku instansi tautologis.}}{}

\tagitem{prvIf}{\iftag{probIf}{Latihan.}{Upamane $!A \lif !B \in
      \Gamma$ lan $!A \in \Gamma$; banjur $\Gamma \Proves[\Sigma] !B$,
      mula $!B \in \Gamma$ amarga katutup sacara deduktif. Kosok baline,
      yen $!A \lif !B \notin \Gamma$, amarga $\Gamma$ $\Sigma$-konsisten
      jangkep, $\lnot (!A \lif !B) \in \Gamma$. Amarga
      $\lnot(!A \lif !B) \lif !A$ iku instansi tautologis, $!A \in
      \Gamma$ amarga katutup sacara deduktif. Amarga $\lnot(!A \lif !B) \lif
      \lnot !B$ iku instansi tautologis, $\lnot !B \in
      \Gamma$. Mula $!B \notin \Gamma$ amarga $\Gamma$
      $\Sigma$-konsisten.}}{}

\tagitem{prvIff}{\iftag{probIff}{Latihan.}{Upamane $!A \liff !B \in
      \Gamma$. Yen $!A \in \Gamma$, mula $!B \in \Gamma$, amarga $(!A
      \liff !B) \lif (!A \lif !B)$ iku instansi tautologis.
      Kanthi cara padha, yen $!B \in \Gamma$, mula $!A \in \Gamma$.
      Dadi $!A \in \Gamma$ lan $!B \in \Gamma$ loro-lorone, utawa
      ora $!A \in \Gamma$ lan ora $!B \in \Gamma$.

Kosok baline, upamane $!A \liff !B \notin \Gamma$. Amarga $\Gamma$
      $\Sigma$-konsisten jangkep, $\lnot (!A \liff !B) \in
      \Gamma$. Amarga $\lnot(!A \liff !B) \lif (!A \lif \lnot !B)$ iku
      instansi tautologis, yen $!A \in \Gamma$, mula $\lnot !B \in
      \Gamma$, lan amarga $\Gamma$ $\Sigma$-konsisten, $!B \notin
      \Gamma$. Kanthi cara padha, yen $!B \in \Gamma$, mula $!A \notin \Gamma$.
      Yen loro-lorone ora dadi anggota himpunan, amarga jangkep,
      negasine loro-lorone dadi anggota. Instansi tautologis
      $\lnot !A \lif (\lnot !B \lif (!A \liff !B))$ lan sifat katutup
      sacara deduktif banjur njalari bikondisional iku dadi anggota,
      sawijining kontradiksi. Dadi ora mungkin $!A \in \Gamma$ lan $!B \in \Gamma$ bebarengan,
      utawa $!A \notin \Gamma$ lan $!B \notin \Gamma$ bebarengan.}}{}
  \end{enumerate}
\end{proof}

\begin{probtag}{probAnd,probOr,probIf,probIff}
  Rampungna bukti \olref[nml][com][ccs]{prop:ccs-properties}.
\end{probtag}

\end{document}
